\documentclass[12pt]{article}
\usepackage[T1]{fontenc}
\usepackage[utf8]{inputenc}
\usepackage{lmodern}
\usepackage{textcomp}
\usepackage{enumitem}
\usepackage{geometry}
\geometry{margin=1in}
\begin{document}
\begin{center}
Answer Key - History - Undergraduate Level Assessment 22 \
\small (Answers are ordered exactly as in the question paper.)
\end{center}

\section*{Multiple Choice}
\begin{enumerate}[label=\arabic*.]
\item \textbf{Answer:} A
\\ \textit{Options:} A) 5100 B.P.; B) 6200 B.P.; C) 7000 B.P.; D) 4000 B.P.; E) 1200 B.P.
\\ \textit{Note:} The date of 5100 B.P. reflects the earliest known archaeological evidence of metal use in South America, indicating the region's advanced technological developments during that time period.
\item \textbf{Answer:} C
\\ \textit{Options:} A) a diet high in calcium or low in calcium.; B) a diet high in saturated fats or low in saturated fats.; C) mostly grains or mostly nuts and fruits.; D) mostly root vegetables or mostly leafy greens.; E) a diet rich in sugar or low in sugar.
\\ \textit{Note:} The correct answer is based on the fact that different plants, such as grains and nuts/fruits, utilize distinct carbon fixation pathways (C$_{3}$ and C$_{4}$), leading to variations in the stable carbon isotope ratios, particularly 13 C, found in the bones of individuals who consumed them.
\item \textbf{Answer:} C
\\ \textit{Options:} A) understanding wave patterns; B) cloud patterns; C) magnetic orientation to the poles; D) bird flight paths; E) currents and wind patterns
\\ \textit{Note:} The native navigators of the Pacific relied on celestial navigation, ocean currents, and environmental cues rather than magnetic orientation to the poles, which was not part of their traditional navigation techniques.
\item \textbf{Answer:} A
\\ \textit{Options:} A) Ming Dynasty Ceramic Art; B) Qin manufacture; C) Han Stoneware; D) Lung-shan; E) Hangzhou Style
\\ \textit{Note:} The correct answer is based on the understanding that the Ming Dynasty was notable for its advancements in construction techniques, including the use of stamped or pounded earth in building durable structures, which reflects the era's architectural innovations.
\item \textbf{Answer:} A
\\ \textit{Options:} A) elite-sponsored feasts in which captives of war were beheaded.; B) public humiliation, forced labor, and consumption of hallucinogenic plants.; C) ritualistic battles, consumption of coca leaves, and bloodletting.; D) communal feasts and the torture of captives to be sacrificed.; E) improved diet, forced labor, and the use of coca leaves.
\\ \textit{Note:} The correct answer is based on the historical practice among the Inca where elite-sponsored feasts served as a ceremonial context for human sacrifices, often involving the beheading of war captives, which highlighted the connection between social hierarchy, ritual significance, and warfare in their culture.
\end{enumerate}

\section*{True / False}
\begin{enumerate}[label=\arabic*.]
\setcounter{enumi}{5}
\item \textbf{Answer:} False
\\ \textit{Options:} A) True; B) False
\\ \textit{Note:} The correct answer is: April 3, 1997
\item \textbf{Answer:} True
\\ \textit{Options:} A) True; B) False
\\ \textit{Note:} According to the passage: raising public awareness about fire and improving safety
\item \textbf{Answer:} False
\\ \textit{Options:} A) True; B) False
\\ \textit{Note:} The correct answer is: Selfridges flagship store
\item \textbf{Answer:} True
\\ \textit{Options:} A) True; B) False
\\ \textit{Note:} According to the passage: death
\item \textbf{Answer:} True
\\ \textit{Options:} A) True; B) False
\\ \textit{Note:} According to the passage: stators
\end{enumerate}

\section*{Long Form Response}
\begin{enumerate}[label=\arabic*.]
\setcounter{enumi}{10}
\item \textbf{Answer:} Hyder Ali
\\ \textit{Note:} Expected answer: Hyder Ali
\item \textbf{Answer:} The Cathar Crusade
\\ \textit{Note:} Expected answer: The Cathar Crusade
\end{enumerate}

\vfill
\noindent\textit{End of Answer Key}
\end{document}